% ---------------------------------------------------------------------------
% Paper 01 -- Obstruction certificates for viability under incomplete observation
% Edition v73 (research-article rewrite of v72; see CHANGELOG_v72_to_v73.md)
% Build: tectonic paper01_obstruction_calculus_v73.tex
% ---------------------------------------------------------------------------
\documentclass[11pt,a4paper]{article}
\usepackage[a4paper,margin=26mm]{geometry}
\usepackage[T1]{fontenc}
\usepackage{lmodern}
\usepackage{microtype}
\usepackage{amsmath,amssymb,amsthm}
\usepackage{graphicx}
\graphicspath{{./}{paper01_obstruction_calculus_v73_src/}}
\usepackage{booktabs,array,tabularx,multirow}
\usepackage[font=small,labelfont=bf]{caption}
\usepackage[round,authoryear]{natbib}
\usepackage{enumitem}
\usepackage[colorlinks=true,linkcolor=blue!50!black,citecolor=blue!50!black,urlcolor=blue!50!black]{hyperref}
\usepackage{etoolbox}
\usepackage{xcolor}
\setlength{\bibsep}{3pt plus 1pt}
\allowdisplaybreaks
\emergencystretch=2em

% ---- theorem environments: one counter per section
\theoremstyle{plain}
\newtheorem{theorem}{Theorem}[section]
\newtheorem{proposition}[theorem]{Proposition}
\newtheorem{corollary}[theorem]{Corollary}
\newtheorem{lemma}[theorem]{Lemma}
\theoremstyle{definition}
\newtheorem{definition}[theorem]{Definition}
\newtheorem{example}[theorem]{Example}
\newtheorem{openproblem}[theorem]{Open Problem}
\theoremstyle{remark}
\newtheorem{remark}[theorem]{Remark}

% ---- macros
\newcommand{\R}{\mathbb{R}}
\newcommand{\Viab}{\mathrm{Viab}}
\newcommand{\RViab}{\mathrm{RViab}}
\newcommand{\ERViab}{\mathrm{ERViab}}
\newcommand{\EViab}{\mathrm{EViab}}
\newcommand{\IRViab}{\mathrm{IRViab}}
\newcommand{\cV}{\mathcal{V}}
\newcommand{\cI}{\mathcal{I}}
\newcommand{\cJ}{\mathcal{J}}
\newcommand{\cR}{\mathcal{R}}
\newcommand{\cA}{\mathcal{A}}
\newcommand{\cW}{\mathcal{W}}
\newcommand{\cY}{\mathcal{Y}}
\newcommand{\cS}{\mathcal{S}}
\newcommand{\Post}{\mathrm{Post}}
\newcommand{\Pre}{\mathrm{Pre}}
\newcommand{\Tobs}{T_{\mathrm{obs}}}
\newcommand{\Atube}{\cA_{\mathrm{tube}}}
\newcommand{\Reach}{\mathrm{Reach}}
\newcommand{\supp}{\mathrm{supp}}
\newcommand{\co}{\operatorname{co}}
\newcommand{\dd}{\,\mathrm{d}}
\newcommand{\hyp}[1]{\textup{(#1)}}

\title{\vspace{-1.2em}Obstruction certificates for viability under incomplete observation\\[0.3em]
{\large A calculus of sufficient conditions for nonviability, with exact finite and polyhedral cases}}
\author{Amin Abaee\\[0.3em]
{\normalsize Independent Researcher}\\
{\normalsize \href{https://orcid.org/0000-0002-0019-1842}{ORCID 0000-0002-0019-1842}\ \ $\cdot$\ \ \href{mailto:amin_abaee@ut.ac.ir}{amin\_abaee@ut.ac.ir}}}
\date{Edition v73 --- October 2026}

\begin{document}
\maketitle
\vspace{-1.5em}
\begin{center}\small\emph{Author preprint, not yet refereed; an earlier edition is deposited at Zenodo and figshare \citep{abaee2026preprint}. The Supplementary Material (Sections S1--S9) carries the complete proofs, worked certificates, the probabilistic counterpart, and the verification record.}\end{center}

\begin{abstract}
\noindent
Under perfect measurement the viability kernel of a constrained control system is characterised by tangency conditions, and under incomplete observation the \emph{sufficiency} side of the theory is served by output-feedback regulation conditions and estimation-tube reductions. The complementary task---certifying that \emph{no} observation-based policy can keep the system within its constraints---has no comparable instrument. This paper develops one: a calculus of \emph{obstruction certificates}, each a checkable witness that a prescribed class of observation-based policies fails because the information structure leaves no room for any policy. The calculus is organised by a single selector principle: an observation-based regulator must choose one response for an entire information state, and viability requires that response to be admissible, safe over the review, and recursively viable for every compatible state. Each certificate proves that one of these response correspondences is empty. The central result, the common-action obstruction, shows that disjoint safe-control sets of states merged by the observation exclude every policy of a held-action class, even when each state is viable under full information; it has a Farkas-pair witness for polyhedral data and a Helly-type sparse witness of at most $m+1$ states for $m$ controls. A delayed-information obstruction turns the timing of the first informative observation into an explicit bound, exactly computable as a linear programme over polytope-declared blind-window classes, and a recourse obstruction certifies the complementary failure of post-observation authority in affine systems. In finite systems the calculus is exact: belief recursion is sound and complete and the common-action test is necessary and sufficient one step ahead. The limits of completeness are made precise---static observation reduces to open-loop robust viability, and no certificate pair that reads only the blind window can be complete. A separate certification layer shows that an exact observation-only safety verdict exists if and only if safe-set membership is constant on the observation fibres. The certificates are instantiated on a symbolic resource model, audited against the exact recursion, and read as design requirements on monitoring timing, coarseness, aggregation and bias. The discrete core is machine-checked in Lean~4.
\end{abstract}

\noindent\textbf{Keywords:} viability kernel; incomplete observation; obstruction certificates; robust viability; output feedback; set-membership estimation; monitoring design; sustainability governance

\smallskip\noindent\textbf{MSC 2020:} 49J53; 93B03; 93C41; 91B76

